% ═══════════════════════════════════════════════════════════
% §1 Introduction
% ═══════════════════════════════════════════════════════════
\section{Introduction}
\label{sec:intro}

The proliferation of large language models (LLMs) has democratized AI capabilities, yet the computational cost of post-training, including supervised fine-tuning (SFT) and reinforcement learning from human feedback (RLHF/GRPO), remains a barrier for individual researchers and small teams. Consumer-grade GPUs, particularly AMD's RDNA3 architecture (RX 7900 XTX, 24GB), offer an accessible alternative to expensive data-center hardware. However, unlike NVIDIA's ecosystem with extensive efficiency studies on A100/H100 clusters~\citep{chowdhery2022palm, hoffmann2022chinchilla}, the training efficiency characteristics of consumer AMD GPUs remain largely unexplored in academic literature.

Existing MFU (Model FLOPS Utilization) studies universally assume dedicated tensor cores (NVIDIA Tensor Cores, AMD Matrix Cores on Instinct). RDNA3 consumer GPUs implement matrix operations via WMMA (Wave Matrix Multiply-Accumulate) microcoded vector instructions rather than dedicated silicon, making the spec-sheet peak (122.8 TFLOPS BF16) fundamentally unreachable in practice. This architectural distinction necessitates a new measurement methodology, our \emph{dual-peak MFU} approach, that separately reports utilization against theoretical and practical peaks.

Unsloth~\citep{unsloth2024} is a widely adopted training framework that recently introduced official AMD support in collaboration with the AMD team, reporting up to 2$\times$ speedups on data-center Instinct GPUs equipped with dedicated Matrix Cores (CDNA architecture). However, consumer RDNA~3 lacks any dedicated AI acceleration hardware and relies entirely on WMMA microcoded vector instructions, making efficiency characteristics observed on CDNA-class hardware non-transferable.

\paragraph{Contributions.} We make the following contributions:
\begin{enumerate}
    \item A \textbf{dual-peak MFU reporting methodology} that separately measures utilization against the spec-sheet theoretical peak and a practical GEMM peak, with quantified bias for architectures lacking dedicated tensor cores.
    \item A \textbf{systematic efficiency characterization} of LLM post-training on consumer RDNA3, including framework comparison (Unsloth vs HF, 3 seeds) and component-stacking gap attribution revealing a 75.7\% efficiency gap.
    \item A \textbf{tuning study} across 10 configurations with mechanistic explanations, yielding an actionable optimal recipe for practitioners.
    \item An \textbf{open-source profiler} incorporating a training validity check (grad\_norm$>$0) to ensure reported throughput reflects genuine parameter updates.
\end{enumerate}
